\documentclass[12pt,a4paper]{ctexart}
\usepackage[margin=2.2cm]{geometry}
\usepackage{amsmath,amssymb}
\usepackage{booktabs}
\usepackage{longtable}
\usepackage{array}
\usepackage{hyperref}
\usepackage{enumitem}
\usepackage{xcolor}
\usepackage{titlesec}
\hypersetup{colorlinks=true,linkcolor=blue,urlcolor=blue}

\title{《新程序相对于 795 程序的优势与结构分析》逐页解释词与答辩备注\\
\large 对应详细版 PPT 的正式说明文档}
\author{}
\date{2026-04-17}

\setlist[itemize]{leftmargin=2em}
\setlist[enumerate]{leftmargin=2.2em}
\titleformat{\section}{\Large\bfseries}{\thesection.}{0.6em}{}
\titleformat{\subsection}{\large\bfseries}{\thesubsection.}{0.6em}{}

\newcommand{\pageentry}[4]{
\subsection{第 #1 页：#2}

\paragraph{本页定位}
#3

\paragraph{本页内容详解}
#4

\paragraph{答辩备注}
\begin{itemize}
\item 回答问题时应优先强调本页所服务的整体主线，即“从原始实验脚本出发，逐步构建统一的混合优化系统”，避免把单页内容割裂为孤立结论。
\item 若评审追问“这一页和最终结果的关系”，建议明确指出本页是在回答“为什么这样设计”或“这一部分在整个系统中承担什么角色”，而不是单独证明某个指标最优。
\end{itemize}
}

\begin{document}
\maketitle
\tableofcontents
\newpage

\section{文档说明}

本文档对应文件：
\begin{itemize}
\item \path{D:/Trae products/flat_top light/report_package_20260331/ppt_new_program_vs_795_detailed.tex}
\item \path{D:/Trae products/flat_top light/report_package_20260331/ppt_new_program_vs_795_detailed.pdf}
\end{itemize}

本文档的用途不是替代 PPT，而是为正式汇报、组会答辩或老师逐页追问场景提供一份更完整的讲解底稿。相较于简短发言提纲，本稿做了两类补充：
\begin{enumerate}[label=\arabic*.]
\item 对每一页给出更正式、可复述的“本页内容详解”；
\item 对每一页补充“答辩备注”，用于提示该页最容易被追问的点，以及建议维持的回答口径。
\end{enumerate}

全文按详细版 PPT 的 38 页顺序组织。需要特别说明的是，当前 PPT 的重点并不是简单比较“新程序”和“旧程序”孰优孰劣，而是说明新程序如何在保留原始 \texttt{795} 实验经验的基础上，引入统一 benchmark、学习型初始化器和物理优化收尾，从而形成一条更可复现、可扩展、可持续迭代的技术主线。

\section{逐页解释词与答辩备注}

\pageentry{1}{标题页}{
本页用于界定汇报主题与讨论范围。
}{
本页的核心作用，是明确本次汇报所讨论的对象并不是单一算法，而是一套已经形成结构化主线的平顶光全息系统。标题中“新程序相对于 795 程序的优势与结构分析”包含两层含义：第一，汇报会说明原始 \texttt{795} 程序在实验层面提供了哪些关键经验；第二，汇报会解释当前新程序如何在这些经验基础上完成结构重组，并进一步引入 Lin2025 风格的学习型初始化与 Bowman 风格的物理优化收尾。换言之，本次汇报的重点不是简单报告实验结果，而是系统阐释“为什么当前程序要这样组织，以及这种组织带来了什么能力提升”。

在正式场合，建议从这一页开始就明确三类技术来源：\texttt{795} 提供实验经验相位与目标构造启发，Lin2025 提供学习型初始化思路，Bowman/hybrid 提供物理一致性的最终优化保障。后续所有页面都围绕这一整合路线展开。
}

\subsection{第 2 页：汇报结构}

\paragraph{本页定位}
本页用于建立听众对整场汇报逻辑顺序的预期。

\paragraph{本页内容详解}
本页给出了全文的组织框架，其作用是在正式汇报开始时为听众建立“从哪里开始、为什么这样展开、最后要落到哪里”的清晰路径。整体结构可以概括为四段。第一段回顾原始 \texttt{795} 程序，目的不是做怀旧式介绍，而是让听众理解当前系统的许多关键设计都直接继承或转译自原始实验脚本中的有效经验。第二段介绍新程序的总体思路与 AutoDL notebook 的版本演进，用来说明当前系统不是一次性构建完成，而是在长期试错与整合中逐步形成。第三段拆解新程序的关键模块，回答“现在的系统到底由哪些部分组成，它们各自承担什么职责”。第四段回到结果与结论，解释这些结构化改动究竟在 benchmark 与工程流程上带来了哪些可验证的收益。

如果时间有限，可以在本页先向听众说明本次汇报的核心问题只有三个：原始 \texttt{795} 到底留下了什么、当前新程序到底重构了什么、这种重构到底有没有带来稳定的系统性提升。这样可以帮助后续每一页都围绕主问题展开，而不至于陷入过细的实现细节。

\paragraph{答辩备注}
\begin{itemize}
\item 如果老师问“为什么先讲 795，再讲新程序”，建议回答：因为新程序并不是凭空设计，而是建立在原始实验脚本提供的目标构造、参考相位与优化经验之上，先交代来源有助于说明后续重构的合理性。
\item 如果被问“为什么还要讲版本演进”，建议回答：因为当前 AutoDL notebook 的主线是从多个版本收敛出来的，很多关键结论不是静态设计，而是通过长期试验确定的。
\end{itemize}

\pageentry{3}{795 程序的定位}{
本页用于界定原始 \texttt{795.py} 的角色与历史价值。
}{
本页要说明的核心事实是，原始 \texttt{795.py} 并不是一个通用优化框架，而是一个深度贴合真实实验系统的工程脚本。它的价值主要体现在三点。第一，它面向的是确定的 SLM 分辨率、确定的光学参数和确定的实验链路，因此目标定义、加权方式和优化流程都与真实装置高度绑定。第二，它能够在既定实验条件下稳定地产生可用的平顶光，这说明其内部包含大量已经被实验验证过的有效经验。第三，它虽然不具备现代软件工程意义上的模块化结构，但恰恰因为和实验系统贴得非常近，所以提供了很多当前新系统仍然要继承的关键先验。

因此，在本页中应避免把 \texttt{795} 描述成“旧而过时的脚本”。更准确的说法是：\texttt{795} 是一套强实验导向、强条件绑定、但经验非常浓缩的工程求解方案。当前新程序的许多结构设计，并不是要推翻它，而是要把其中有效但隐含的经验显式化、模块化，并纳入统一 benchmark 体系中。
}

\pageentry{4}{795 的整体运行顺序}{
本页用于解释原始 \texttt{795} 的工作流程为何有效，以及其结构特点是什么。
}{
本页展示了原始 \texttt{795} 程序的基本执行顺序：读取实验参数、生成入射包络、构造目标振幅和目标相位、构造加权掩模、生成初始相位、执行 CG 优化，并在接入实验条件时结合相机反馈再次修正。这个顺序非常重要，因为它说明 \texttt{795} 并不是“先训练一个模型再看效果”的路线，而是一种直接围绕实验对象进行数值求解和反馈修正的流程。

从工程角度看，这种结构的优点非常直接：整条链路短、控制明确、和实验数据贴合紧密，因此一旦针对目标任务调通，就往往有较高的可用性。但其局限也同样明显，即整条流程高度耦合在同一脚本中，缺少统一抽象，很难方便地比较不同初始化路线、不同损失组合或不同 teacher 策略。这也是当前新程序在组织方式上进行重构的直接动因。
}

\pageentry{5}{795 的目标构造}{
本页用于说明原始实验任务与当前通用 flat-top 目标并不完全相同。
}{
本页需要强调的一点是，原始 \texttt{795} 所面对的目标，不是简单的圆形平顶，而更接近于线形平顶或各向异性平顶。其目标构造方式体现为：一个方向上具有有限平顶宽度，另一个方向上保持高斯型衰减。这种目标设计和真实实验需求强绑定，因此不能简单地用当前默认圆形 flat-top 目标去完全等价替代。

这一点之所以重要，是因为后续所有训练、benchmark 和初始化学习都会受到目标定义影响。如果目标本身发生了形状层面的变化，那么 teacher 结构、优化方向和最终评价指标都会随之变化。因此，本页不仅是在介绍原始脚本做了什么，也是在提示听众：当我们说“新程序相对于 795 更好”时，必须明确比较是在什么目标定义下成立的。后来 line-flat 主线的形成，本质上就是对原始目标经验的进一步回归和结构化继承。
}

\pageentry{6}{795 的加权区域设计}{
本页用于解释 \texttt{795} 稳定性的一个关键来源，即有效区域优先优化。
}{
本页的主要内容，是说明原始 \texttt{795} 并没有试图让整个输出面都达到同样的优化标准，而是通过 \texttt{Wcg} 和 \texttt{W99} 等加权掩模，将优化资源集中到真正重要的区域。这一策略带有明显的实验工程特征：实验通常不会同等关注整个输出面，而更关心目标区域内的均匀性、相位平坦性和可用能量。

从方法论上看，这说明原始 \texttt{795} 已经包含了“局部关键区域优先”的思想，只是当时以经验参数和加权掩模的形式存在。新程序中后来出现的 core region、ROI、core-phase-first 排序规则，以及对不同 route 的统一 benchmark，实际上都是在把这种工程经验显式化和标准化。因此，本页是旧程序与新程序之间的一座关键桥梁。
}

\pageentry{7}{795 的初始相位与传播模型}{
本页用于说明原始 \texttt{795} 的初值设计具备明确的物理先验含义。
}{
本页解释的是 \texttt{795} 初始相位为何并非普通意义上的随机初值。原始脚本中的 \texttt{phase\_guess} 通常包含线性倾斜项、二次曲率项以及整体分布调制项，这些成分分别对应输出偏移、包络尺度与整体形状控制。因此，初始相位本身就是一种对实验系统条件的编码，而不是一个与任务无关的数值起点。

这一认识对于当前系统非常重要。新程序中的 \texttt{reference} 路线并不是机械地“读取一张旧相位图”，而是在继承原始经验型初值的思想基础上，将其转化为可解析、可插值、可适配当前目标定义和仿真分辨率的初始化来源。也就是说，本页帮助说明为什么当前系统会高度重视 \texttt{795} 参考相位，以及它为何能够作为学习型初始化器之外的一条重要路线长期保留。
}

\pageentry{8}{795 的优化方式}{
本页用于说明原始 \texttt{795} 已经具备多目标优化的雏形。
}{
本页的重点，是解释原始 \texttt{795} 并非只依赖单一 cost function，而是在不同阶段使用了多类优化目标，例如加权区域内的振幅误差、基于 overlap 的保真度项，以及 overlap 与效率惩罚的组合。这说明原始脚本虽然没有形成今天这种统一 benchmark 语言，但实际上已经在工程层面处理“多个质量指标之间如何平衡”的问题。

这为当前系统的损失设计与 benchmark 体系提供了历史依据。当前新程序中将核心区均匀性、核心区相位平坦度和效率显式列为三项主要指标，并通过统一脚本比较不同路线，本质上是在对原始 \texttt{795} 中隐式存在的多目标平衡思想进行重新组织和结构化表达。
}

\pageentry{9}{795 的实验闭环}{
本页用于强调原始 \texttt{795} 的真实实验价值。
}{
本页说明原始 \texttt{795} 的一个关键能力，即它能够接入真实相机，基于测量结果修正目标，再重新执行优化，最终将相位回写到 SLM。这种“测量 - 修正 - 再优化”的真实闭环，是许多纯离线优化流程所不具备的。也正因为如此，\texttt{795} 的优势不能只用仿真输出的一组静态指标来完全概括。

在讲解这一页时，建议明确区分“当前新程序已经形成统一仿真与训练主线”和“当前新程序尚未完全恢复原始实验闭环”这两件事。前者说明系统主线已经建立，后者则说明它与原始实验控制版本之间仍然存在明确差距。这种区分有助于避免听众误以为当前系统已经在所有维度上完全替代 \texttt{795}。
}

\pageentry{10}{795 的优点与局限}{
本页用于对原始 \texttt{795} 做平衡性的阶段总结。
}{
本页的任务是给出一个公允的评价。原始 \texttt{795} 的主要优点在于：最贴近真实实验，目标与加权设计经验非常充分，参考相位具备直接实验价值，并且能够接入相机反馈闭环。其主要局限则在于：对具体实验系统的依赖非常强，泛化能力有限，不便于统一比较多条路线，也难以自然融入现代训练流程和远程迭代工作流。

因此，本页之后应顺势引出新程序的设计动机：当前系统并不是为了“否定 795”，而是试图解决 \texttt{795} 在结构化、可扩展与可比较性方面的不足，同时尽可能保留其对实验最有价值的经验部分。
}

\pageentry{11}{Bowman：参考路线的定位}{
本页用于把 Bowman 从“一个名字”讲成一条明确的方法路线。
}{
本页的核心任务，是说明 Bowman 不应被狭义理解为“某一个 CG 代码实现”，而应被理解为一类直接面向物理目标函数求解的思路。它显式保留传播模型、显式定义输出场上的目标函数，并直接在 SLM phase 上进行数值优化。因此，Bowman 的真正价值在于它的物理可解释性和收尾稳定性，而不只是它是否使用了共轭梯度这一具体优化器。

在这一页中，建议把 Bowman 的优势和局限同时讲清楚。它的优势是：目标、传播与优化之间关系清楚，结果具备较强物理一致性；它的局限则是：如果初值较差，直接从头搜索往往代价较高，且对初始化条件较为敏感。这个结论正好为后面“为什么需要 Lin”做铺垫。
}

\pageentry{12}{Bowman：源码中的落点与当前吸收方式}{
本页用于解释 Bowman 在当前系统里不是概念性引用，而是有清楚的代码落点。
}{
本页建议从“思想来源”和“代码落点”两方面展开。思想来源方面，当前项目中对 Bowman 的理解并不只来自二手概述，而是结合了 \path{references/Bowman/} 中的原始文献，以及 \path{references/ftl_gen/795.py} 里保留下来的实验目标构造与优化经验。代码落点方面，则主要体现在 \path{ai_holography/hybrid.py} 的 \texttt{bowman\_cg\_refine(...)} 和 \path{ai_holography/scripts/compare_bowman_ai.py} 这类对比入口中。

更重要的是，本页需要讲清楚当前系统并没有把 Bowman 原封不动复制过来，而是进行了现代化重组。现在的做法是，把它保留为“最终物理收尾机制”，并通过 \texttt{efficiency\_overlap}、\texttt{composite} 和 \texttt{phase\_priority} 等分阶段目标逐步收束。也就是说，Bowman 在当前系统中承担的是“最后一公里”的职责，而不再独自承担从任意初值起步的全部负担。
}

\pageentry{13}{Lin：参考路线的定位}{
本页用于把 Lin 从“训练一个网络”提升为一条有明确问题改写思路的方法路线。
}{
本页最需要讲清楚的一点是，Lin 路线的核心并不是“让 CNN 直接输出最终 hologram”，而是把原问题改写到更适合学习的表示域中。具体来说，它倾向于让网络先在位置域预测 amplitude 和 phase，再通过相应变换恢复 hologram 初值。这种做法背后的目的，是避免直接在最终相位空间中进行过于刚性的回归，从而提升训练稳定性，并让网络更容易学习到与目标结构相关的统计规律。

因此，在这页里可以把 Lin 的核心价值概括为“学习初始化规律”。它的优势是：能够从数据中学习出什么样的初值更可能成功；它的局限则是：如果单独把它当作最终求解器使用，往往还不足以完全替代后端物理优化。这一结论也为后面的混合路线说明提供了逻辑基础。
}

\pageentry{14}{Lin：源码中的落点与当前吸收方式}{
本页用于说明当前系统对 Lin 的吸收不是表面复刻，而是面向平顶光任务的再设计。
}{
本页可以直接对应 \path{lin2025_holography/README.md}、\path{lin2025_holography/model.py} 和 \path{lin2025_holography/training.py}。其中，\texttt{Lin2025HologramNet} 明确采用位置域 amplitude / phase 预测，\texttt{position\_to\_hologram(...)} 负责把预测结果变回 hologram 初值；而训练循环则不仅包含基本监督项，还叠加了 \texttt{teacher phase distillation} 和面向 flat-top 任务的质量损失。

这说明当前系统对 Lin 的吸收，并不是机械复刻一篇论文结构，而是明确将其改造成“学习型初始化器”模块。评价一个 Lin checkpoint 的好坏，也不再只看 supervised loss，而要看它作为 init 时，能否让后续 hybrid polish 更容易进入更优 basin。换言之，Lin 在这里服务的是全链路最终质量，而不是局部监督误差最小化。
}

\pageentry{15}{为什么当前系统最终采用 Bowman + Lin 的混合路线}{
本页用于给出“为什么不是二选一，而是混合”的直接回答。
}{
这页的核心逻辑非常适合在答辩中直接复述。如果只有 Bowman，系统会有很强的物理一致性和结果稳定性，但容易受初值影响；如果只有 Lin，系统会获得快速、可学习的初始化能力，但单独作为最终求解器时仍然不够稳。因此，当前混合系统选择让 Lin 先解决“从哪里开始更好”，再让 Bowman 解决“最后怎样收得更稳、更符合物理约束”。

从系统结构上看，这种混合不是简单串联，而是由 runner 负责 route 选优，由 benchmark 负责结果验证。也就是说，Lin 和 Bowman 在当前系统里并非竞争关系，而是前后分工关系。这个结论非常关键，因为它解释了为什么当前系统的创新点不是抛弃传统方法，而是在其基础上完成结构增强。
}

\pageentry{16}{四条路线的角色对照：795 / Bowman / Lin / 当前混合系统}{
本页用于把前面几条路线的角色、强弱项和在当前系统中的位置压缩成一张总表。
}{
这一页的作用是“收口”。前面几页分别讲了 \texttt{795}、Bowman 和 Lin 的思想来源与代码落点，而这一页则把四条路线放进同一张表中，明确比较它们的核心定位、主要强项、主要短板、当前在代码中的对应位置，以及在答辩中最适合采用的表述方式。

这类总表在正式汇报中特别有用，因为它能够帮助听众快速把前面的信息组织起来，也方便在老师追问“那你到底是在做 795、Bowman、Lin 还是别的东西”时，给出一句明确回答。最推荐的口径是：\texttt{795} 是历史实验先验与参考基线，Bowman 是最终物理收尾机制，Lin 是学习型初始化器，而当前混合系统才是本次真正要交付和解释的主线。
}

\pageentry{17}{新程序的总体设计目标}{
本页用于明确新程序的最高层设计原则。
}{
本页要传达的中心思想是：当前新程序不是某一种单独方法的替代，而是三类路线的整合。其一，保留 \texttt{795} 所提供的实验先验和参考相位资源；其二，保留 Bowman/hybrid 的物理一致性和最终收尾能力；其三，引入 Lin2025 风格的学习型初始化器，以缓解纯物理优化对初值敏感的问题。

因此，新程序的目标不只是追求某个单项指标更低，而是构建一个更完整的系统，使训练、推理、benchmark、远程导出和本地分析能够统一起来。换言之，新程序的核心价值是“系统能力的形成”，而不仅仅是算法单点上的替换。
}

\pageentry{18}{新程序整体运行流程}{
本页用于说明当前系统的主线是如何从目标定义一直贯通到 benchmark 输出的。
}{
本页给出了一条完整的系统流程：先构造目标场和参考信息，再生成多个初始化候选，之后通过候选选优得到更优初值，然后进行 phase-only correction、多分辨率 refine 和 Bowman-style hybrid polish，最后再输出统一 benchmark 和分析结果。这条流程说明当前系统是一个串联的、角色明确的混合框架，而不是将多种方法简单并列堆叠。

在正式讲解时，应特别强调神经网络在当前系统中的定位。它并不负责直接输出最终实验可用的最优相位，而是作为学习型初始化器，为后续物理优化创造更好的起点。最终结果的质量仍然依赖物理收尾阶段来保证。这一点有助于避免听众把当前系统误解为“纯神经网络替代传统物理算法”的路线。
}

\pageentry{19}{新程序的目录结构}{
本页用于说明当前工程拆分背后的设计逻辑。
}{
本页展示了当前工程的主要目录分层。原始实验资产保留在 \texttt{references/ftl\_gen} 中，用于保存 \texttt{795} 脚本、参考相位和历史 CG 模块；\texttt{ai\_holography} 是当前物理主线代码所在，负责目标构造、传播、初始化选优与 hybrid polish；\texttt{lin2025\_holography} 则承担数据集、网络、训练和预测相关职责；\texttt{colab\_local\_split} 负责远程训练和本地分析工作流；\texttt{report\_package\_20260331} 则负责汇报整理和结果归档。

这一结构拆分的意义在于，原来耦合在一个实验脚本里的内容，如今已被分为具有清晰边界的若干模块。这样不仅有利于维护，也有利于明确回答“某个结论来源于哪一层设计”以及“后续优化应该从哪一层入手”的问题。
}

\pageentry{20}{版本演进总览}{
本页用于总览 AutoDL notebook 从早期到当前稳定交付版的演进脉络。
}{
本页首先将 notebook 的版本演进划分为六个阶段：从最初的自包含工作流搭建期，到 legacy task 与 flat-top 调参期，再到 line-flat 主线形成、集成 notebook、以及高分辨率和结构一致性修复阶段。其核心作用并不是列举版本号，而是帮助听众理解当前系统并不是一次性设计完成，而是在长期迭代中逐步收敛到今天的形态。

本页还特别指出了若干已知元数据漂移问题，例如标题与文件版本不一致、说明文字沿用旧版本等。这一说明很重要，因为它解释了为什么版本史需要根据 notebook 的实际内容和关键结构变化来重建，而不能只依赖文件名或首页标题。当前最值得记住的关键拐点是 \texttt{v30}、\texttt{v42a}、\texttt{v43}、\texttt{v53} 和 \texttt{v56}，后续若被追问版本演进，建议始终围绕这几个节点组织回答。
}

\pageentry{21}{v2-v18 / v19-v39：从可运行到主线成形}{
本页用于解释早中期版本如何从“能跑”逐步走向“形成主线”。
}{
本页的前半部分对应 \texttt{v2-v18}，这一阶段主要解决的是基础性问题，如 notebook 自包含结构、远程训练可运行性、teacher 池构造、OOM 回退策略，以及 line-flat 任务定义与数据组织。换句话说，这一时期更像是在搭建实验平台和训练框架，而不是追求最终最优结果。

后半部分对应 \texttt{v19-v39}，其核心在于逐步验证 learned init 与 hybrid 结合这条路线的可行性，并最终在 \texttt{v30} 附近形成第一代真正可复用的主线。随后若干版本虽然继续做了质量分支扫参和局部探索，但最终结论是：\texttt{v30} 的主方向是正确的，后续工作更应该围绕它收束，而不是无休止扩展旁支路线。
}

\pageentry{22}{v40-v50：集成 notebook 与历史表阶段}{
本页用于解释中期版本为什么要强调“单变量验证”和“收束”。
}{
这一阶段的主要任务，不是继续扩充方法家族，而是把此前探索中真正有效的改动识别出来，并重新集成为更稳定的基座。例如，\texttt{v42a} 只改参考相位生成中的 \texttt{efficiency\_floor}，用于验证该变量本身的影响；\texttt{v42b} 只改 ROI 覆盖，以判断该因素是否足以构成默认主线；\texttt{v43} 则将 checkpoint 路径从共享目录拆开，避免代码版本与权重版本混淆。最终，这些有效改动被整合进 \texttt{v44}，形成更加干净的统一底座。

这一页的关键结论是：中期版本真正重要的，不是尝试了多少路，而是建立了“哪些变化值得留下、哪些变化只是过渡尝试”的判断标准。这种判断能力对于后续高分辨率和结构一致性修复阶段尤为重要。
}

\pageentry{23}{v51-v58：高分辨率、结构一致性与当前推荐版本}{
本页用于解释为什么当前推荐交付的是 \texttt{v58}，以及它与关键算法拐点之间的区别。
}{
本页展示了后期版本演进的主线。首先，\texttt{v52} 将高分辨率和 line-flat 目标真正拉入主流程，但也因此暴露出效率塌陷问题；随后 \texttt{v53} 通过修复 \texttt{efficiency\_floor} 和相关权重，使高分辨率 line-flat 主线重新回到可用区间，这一步构成后期最重要的算法修复拐点之一。再往后，\texttt{v56} 开始处理 notebook 结构一致性问题，例如统一 floor、清理重复逻辑、保证输出设置一致；\texttt{v58} 则在此基础上补回导出代码、统一 summary 路径、补齐说明字段，并修正运行目录与版本一致性。

因此，当前推荐交付版本是 \texttt{v58}，因为它最适合作为稳定附件与复现实验入口；但从算法主线角度看，真正决定方向的关键节点仍然是 \texttt{v30 / v42a / v43 / v53 / v56}。这是本页最需要明确传达的区分。
}

\pageentry{24}{模块 1：目标与参考相位模块}{
本页用于解释目标构造和参考相位管理为何是当前系统的基础层。
}{
该模块主要由 \texttt{targets.py} 和 \texttt{references.py} 组成，其作用可以概括为两点。第一，统一构造当前任务所使用的目标振幅、目标相位和相关区域定义，使后续训练、推理与 benchmark 都建立在一致的问题描述之上。第二，将历史 \texttt{795\_*.npy} 参考相位转化为当前系统中可解析、可插值、可适配的工程资产。

这一层的价值在于把原本隐含在实验脚本中的经验显式化。过去很多“有效但说不清”的参考相位，现在都能以结构化方式参与 route 比较与初始化选优。因此，本页实际上是在说明：当前系统为何能够既继承旧经验，又不被旧脚本结构所束缚。
}

\pageentry{25}{模块 2：传播与损失模块}{
本页用于说明整套系统共享的物理和评价底盘。
}{
传播模块与损失模块共同定义了当前系统的“公共物理语言”。传播模块保证从 SLM 面到输出面的建模具有明确的物理意义；损失模块则负责对 overlap、核心区均匀性、核心区相位平坦度和效率等关键指标进行统一定义。由于这些指标在训练阶段、refine 阶段和 hybrid polish 阶段都会被调用，因此它们不仅是最终结果评估标准，也是整个系统各条路线得以公平比较的前提。

在答辩中，如果需要说明“为什么当前 benchmark 有说服力”，可以回到本页。因为 benchmark 的价值并不只是把若干路线放在一张表里，而是这些路线从目标定义到损失计算都共享同一套物理和评价框架。
}

\pageentry{26}{模块 3：AIHolographyPipeline}{
本页用于说明系统编排层的作用。
}{
\texttt{AIHolographyPipeline} 的主要职责是将原本分散的目标构造、初始化预测和多分辨率 refine 组织为一条可复用的流程。它既承接了传统物理传播和优化主线，也为插入 AI 初始化器和不同任务 profile 提供了统一接口。因此，它更像是系统中的“流程编排层”，而不是单独的算法模块。

这类设计的工程意义在于，一旦后续需要增加新的初始化器、校正项或不同目标族，原则上不需要再从头堆叠新的 notebook 逻辑，而可以在现有 pipeline 上完成扩展。也就是说，本页体现的是系统从“脚本拼接”向“模块编排”转变的关键一步。
}

\pageentry{27}{模块 4：HybridHolographyRunner}{
本页用于解释候选选优和 route 调度机制为何是当前系统的重要结构改进。
}{
\texttt{HybridHolographyRunner} 负责统一管理不同来源的初始化候选，包括 \texttt{reference}、\texttt{lin2025} 和 \texttt{warm\_start} 等。它会先生成候选，再通过 overlap 和初始化评分机制比较它们的质量，选出更优起点后继续执行 correction、refine 和 hybrid polish。与原始流程相比，这一设计将“初值敏感”从一个隐性问题，转化为一个被显式建模、可重复比较的流程环节。

本页的核心意义在于，它解释了为什么当前系统能够在同一套框架下比较不同初始化路线，并从中挑选更适合当前任务的起点。换言之，新程序在这里做的不只是“多放几个候选”，而是建立了一套统一的 route 管理与选优机制。
}

\pageentry{28}{模块 5：Bowman-style hybrid polish}{
本页用于强调物理优化在当前系统中的最终保障作用。
}{
这一模块直接承接 Bowman 风格的共轭梯度物理优化思想，其内部通过多个阶段逐步在效率、overlap、均匀性和相位平坦性之间完成收束。当前系统之所以能在引入 AI 初始化的同时保持结果可解释、可控，很大程度上依赖于这一模块作为最终物理收尾。

因此，在汇报时应明确指出：新程序并不是“用 AI 取代 Bowman”，而是“用 AI 改善 Bowman 最容易受影响的初值问题，再由 Bowman 风格的物理优化负责最终结果”。这样既能体现创新点，也能避免让评审误解为系统对物理一致性的重视被削弱了。
}

\pageentry{29}{模块 6：Lin2025 数据集与 teacher 策略}{
本页用于说明学习型初始化器是如何吸收实验经验和物理优化经验的。
}{
当前数据集设计并非只依赖单一 teacher，而是同时利用 \texttt{795} 参考相位与系统内部已获得的高质量 hybrid 结果作为监督来源。这意味着网络在训练过程中既能继承原始实验中有效的结构先验，也能吸收当前系统中经过物理优化验证的高质量解特征。teacher 由此不再是一个静态资源，而成为可以随着 benchmark 结果持续更新的“知识池”。

这种设计的重要意义在于，网络学习到的不只是“拟合一张 teacher 图”，而是逐渐形成“什么样的初始化更有利于后续物理收尾”的经验。因此，本页是理解 Lin2025 在当前系统中角色定位的关键页面。
}

\pageentry{30}{模块 7：Lin2025 网络结构}{
本页用于解释网络为什么采用位置域输出而不是直接回归 hologram 相位。
}{
当前网络结构的核心思想，是在位置域预测 amplitude 与 phase，再通过相应变换得到 hologram 初值。这种做法继承了 Lin2025 的基本思路，即避免直接在频域或最终相位空间中进行过于刚性的回归。通常来说，这会带来更稳定的训练过程，也更方便和后续物理传播模型衔接。

本页在答辩中的主要作用，是回应“为什么不直接让网络输出最终相位”这一常见问题。当前系统之所以不这样做，是因为最终结果仍需要依赖物理收尾来保证稳定性，而网络最适合承担的是初始化器角色，而不是终局求解器角色。
}

\pageentry{31}{模块 8：Lin2025 训练循环}{
本页用于说明训练目标是多重的，而不是单一监督拟合。
}{
训练主循环中包含多个损失分量：振幅监督、相位监督、hologram 场误差、teacher phase distillation，以及专门面向 flat-top 质量的损失。这说明训练目标并不局限于“让网络尽量复现 teacher”，而是希望网络输出一种更适合后续物理优化继续改进的初始相位。

从工程角度看，这样的复合损失设计有两个好处。第一，它让训练目标更贴近最终任务质量，而不是只停留在中间表征层面的拟合。第二，它允许在训练过程中输出多个不同侧重的 checkpoint，例如 \texttt{best\_supervised}、\texttt{best\_quality} 和 \texttt{best\_hybrid\_polish}，从而为后续 benchmark 留出更大的筛选与分析空间。
}

\pageentry{32}{模块 9：统一 benchmark}{
本页用于说明为什么 benchmark 是当前系统形成有效结论的基础设施。
}{
统一 benchmark 的作用，是将 \texttt{795\_only}、\texttt{hybrid\_only}、\texttt{lin2025\_only}、\texttt{lin2025\_plus\_hybrid\_balanced} 和 \texttt{lin2025\_plus\_hybrid\_quality} 等多条路线放进同一评价口径中，使用统一指标进行比较。这意味着所有结果不再停留在“看起来更好”或“感觉更稳”这种经验判断上，而能够通过结构一致的方式被明确对照。

本页尤其适合回答“当前新程序到底比 795 好在哪里”这一问题。因为只有在统一 benchmark 存在的前提下，系统性结论才具备可复现性和说服力。从这个意义上说，benchmark 不是结果展示的附属工具，而是当前整个系统得以持续演进的核心基础设施之一。
}

\pageentry{33}{模块 10：远程训练与本地分析}{
本页用于说明当前工作流层面的工程改进。
}{
当前工作流采取“远程平台训练并导出 checkpoint bundle，本地统一分析并与官方 best 对比”的组织方式。这种设计的直接好处是，将 GPU 训练资源的使用与本地统一 benchmark 口径解耦，从而使模型训练、checkpoint 管理和结果分析各自有更加稳定的运行环境。当前推荐的远程 notebook 是 \texttt{train\_lin2025\_autodl\_lineflat\_highres\_export\_v58.ipynb}，其价值主要在于结构、导出和 summary 路径已较为稳定。

本页在正式汇报中的作用，是说明当前系统不仅有算法主线，也已经形成较完整的工程工作流。这一点对于答辩或项目汇报很重要，因为它表明系统具备持续迭代能力，而不是只能在单一环境中偶然跑通一次。
}

\pageentry{34}{结果对比}{
本页用于概括当前系统在统一 benchmark 下的总体表现。
}{
本页的表格展示了多条主要路线在核心区均匀性损失、核心区相位平坦度和效率上的对比结果。讲解这一页时不建议逐个数字机械朗读，而应重点说明整体趋势。首先，相对于 \texttt{795\_only}，当前新程序在核心区相位平坦性和整体可控性上已有显著改善。其次，相对于 \texttt{hybrid\_only}，引入学习型初始化之后，前两项关键质量指标往往还能进一步下降。再次，远程训练得到的最佳 checkpoint 说明，这套系统仍然具备继续向更优结果推进的空间。

因此，本页最重要的结论并不是某一个单独数字，而是：当前混合系统已经具备了稳定产生有效改进的能力，并且这种改进可以通过统一 benchmark 方式被明确展示出来。
}

\pageentry{35}{相对于 795 的优势}{
本页用于总结当前系统相对于原始 \texttt{795} 的系统性提升。
}{
相对于原始 \texttt{795}，当前新程序的提升主要体现在四个方面。第一，不再依赖单一经验参考相位，而是允许多种初始化来源并行存在并进行选优。第二，建立了统一 benchmark，使不同路线可以在相同口径下被比较。第三，引入了学习型初始化器，从而使系统具备通过训练持续提升初值质量的能力。第四，远程训练与本地分析工作流的形成，使 checkpoint 的管理、对比与迭代更加规范。

需要强调的是，这里的优势不应只被理解为“某个指标更低了”，而应理解为“方法论与工程组织方式更成熟了”。也就是说，当前系统相对于 \texttt{795} 的重要进步，是从单一路线的经验求解，发展为可对比、可迭代、可复现的混合优化系统。
}

\pageentry{36}{相对于 Bowman/hybrid 的价值}{
本页用于说明当前系统与纯物理优化路线之间的关系。
}{
本页要表达的重点是，当前系统并不是要取代 Bowman/hybrid，而是要补强其最大的短板，即对初值的敏感性。通过引入 Lin2025 风格的学习型初始化器，后端物理优化更容易从较好的 basin 出发，从而在部分结果上找到优于纯物理基线的 Pareto 点。与此同时，最终物理一致性和可解释性仍然由 Bowman-style hybrid polish 负责保障。

因此，本页最合适的表述是“补强而非替代”。这一点在答辩场合尤其重要，因为它有助于避免将当前工作误解为简单的深度学习替代传统方法，而是更准确地呈现为“物理主线 + 学习型初始化”的混合增强方案。
}

\pageentry{37}{当前不足}{
本页用于主动说明当前系统的边界与待解决问题。
}{
当前系统虽然已经形成较完整的训练、初始化与物理优化主线，但仍存在若干尚未完全解决的问题。第一，默认目标与原始实验中的线形平顶尚未完全一致，因此任务定义仍有继续对齐空间。第二，默认 notebook 流程仍以仿真闭环为主，尚未把真实相机测量重新整合进默认主线。第三，当前训练与实验分辨率仍未完全统一，因此高分辨率迁移和实验对齐仍是后续重要工作。第四，checkpoint 之间仍有一定波动，说明模型稳定性和筛选机制还有继续加强的必要。

在正式汇报中，应将这些问题表述为“系统已知边界”和“下一步工作方向”，而不是简单归类为缺陷。这种表述既诚实，也更能体现对当前系统状态的清晰判断。
}

\pageentry{38}{最终结论}{
本页用于收束全文，并给出当前阶段的总判断。
}{
本页的中心结论是：当前新程序最重要的改进，并不是简单替换某一类算法，而是在工程上完成了 \texttt{795} 实验先验、Lin2025 学习型初始化和 Bowman 物理优化主线的系统整合。系统当前已经具备统一 benchmark、远程训练、本地分析和 checkpoint 选优等完整能力，因此其价值体现在“主线已经成形”，而不只是“某一次实验成功”。

同时，本页还明确指出当前推荐交付的 notebook 是 \texttt{v58}，而真正决定当前算法方向的关键版本是 \texttt{v30 / v42a / v43 / v53 / v56}。这种区分有助于听众理解：交付稳定版和算法关键拐点并不总是同一个版本。最后的收束建议是，将当前阶段概括为“方法有效性已基本确认，后续重点将转向一致性、复现性和实验对齐程度的继续提升”。
}

\section{逐页高频追问与建议回答}

\subsection{第 1 页：标题页}
\begin{itemize}
\item \textbf{问题：为什么标题强调“结构分析”，而不只强调结果提升？}

\textbf{建议回答：} 因为当前工作的核心价值不只是一组更优指标，而是已经形成了由实验先验、学习型初始化和物理优化组成的完整系统。只有先把结构讲清楚，后面的结果才有可解释性。

\item \textbf{问题：这次汇报最想证明的一件事是什么？}

\textbf{建议回答：} 最想证明的是，当前系统已经不再是零散试验集合，而是一条可训练、可比较、可复现的混合优化主线。
\end{itemize}

\subsection{第 2 页：汇报结构}
\begin{itemize}
\item \textbf{问题：为什么要先讲旧程序，再讲新程序？}

\textbf{建议回答：} 因为新程序很多关键设计直接继承自原始 \texttt{795} 的实验经验。先交代来源，后面的重构逻辑才会更清楚。

\item \textbf{问题：为什么版本演进要单独占一部分？}

\textbf{建议回答：} 因为当前 notebook 主线是从长期迭代中收敛出来的，许多关键结论不是静态设计，而是逐版验证得到的。
\end{itemize}

\subsection{第 3 页：795 程序的定位}
\begin{itemize}
\item \textbf{问题：原始 \texttt{795} 到底是算法，还是工程脚本？}

\textbf{建议回答：} 更准确地说，它是围绕真实实验系统构建的工程优化脚本，其中包含有效算法思想，但整体组织方式更偏实验工程而非通用框架。

\item \textbf{问题：为什么不能直接继续沿用 \texttt{795}？}

\textbf{建议回答：} 因为它和具体实验条件绑定过强，不便于统一比较不同路线，也不利于接入现代训练与远程迭代工作流。
\end{itemize}

\subsection{第 4 页：795 的整体运行顺序}
\begin{itemize}
\item \textbf{问题：原始 \texttt{795} 的最大优势是什么？}

\textbf{建议回答：} 最大优势是链路短、实验针对性强，一旦参数调通，往往能稳定地产生真实实验可用结果。

\item \textbf{问题：它的最大局限是什么？}

\textbf{建议回答：} 最大局限是结构耦合严重，难以模块化复用，也不便于在统一框架下比较多条初始化或优化路线。
\end{itemize}

\subsection{第 5 页：795 的目标构造}
\begin{itemize}
\item \textbf{问题：为什么要特别强调线形平顶和圆形平顶的差别？}

\textbf{建议回答：} 因为目标定义会直接影响 teacher、训练方向和最终 benchmark 结论。如果目标本身不同，那么简单比较结果大小就可能失去严格意义。

\item \textbf{问题：当前新程序已经完全恢复到原始线形平顶了吗？}

\textbf{建议回答：} 还没有完全等价恢复，但后期主线已经明显向 line-flat 方向回收，这也是 \texttt{v52-v58} 阶段的重要内容。
\end{itemize}

\subsection{第 6 页：795 的加权区域设计}
\begin{itemize}
\item \textbf{问题：为什么 \texttt{795} 的区域加权思想值得保留？}

\textbf{建议回答：} 因为实验中真正重要的通常不是整个输出面，而是目标有效区域。把优化资源集中到关键区域，是它稳定性的核心来源之一。

\item \textbf{问题：当前系统中的 core region 和 ROI 与这里是什么关系？}

\textbf{建议回答：} 可以理解为对原始区域加权思想的显式化和标准化表达，使其能进入统一 benchmark 和 route 比较。
\end{itemize}

\subsection{第 7 页：795 的初始相位与传播模型}
\begin{itemize}
\item \textbf{问题：为什么说 \texttt{795} 初相位不是“普通初值”？}

\textbf{建议回答：} 因为其中包含了偏移、尺度和整体分布等物理先验，实质上是把实验经验编码进了初始相位。

\item \textbf{问题：当前的 \texttt{reference} 路线和原始 \texttt{795} 的关系是什么？}

\textbf{建议回答：} 当前 \texttt{reference} 路线是对原始经验型初值的结构化继承和适配，而不是简单原样照搬。
\end{itemize}

\subsection{第 8 页：795 的优化方式}
\begin{itemize}
\item \textbf{问题：原始 \texttt{795} 有没有多目标优化思想？}

\textbf{建议回答：} 有，只是它以工程经验形式存在，没有今天这种统一 benchmark 语言。它已经在振幅、overlap 和效率之间做权衡。

\item \textbf{问题：当前损失设计和原始 \texttt{795} 是断裂关系吗？}

\textbf{建议回答：} 不是。当前很多指标化表达，本质上是在把原始脚本中的经验目标重新组织成统一的评价体系。
\end{itemize}

\subsection{第 9 页：795 的实验闭环}
\begin{itemize}
\item \textbf{问题：当前新程序已经接回真实相机了吗？}

\textbf{建议回答：} 默认主流程还没有完全接回真实相机闭环，目前主要仍是仿真闭环，但代码结构上已经预留了相应接口。

\item \textbf{问题：为什么这件事要特别强调？}

\textbf{建议回答：} 因为这决定了当前系统和原始实验控制流程之间还有多大距离，也关系到后续是否能平滑迁移回真实实验环境。
\end{itemize}

\subsection{第 10 页：795 的优点与局限}
\begin{itemize}
\item \textbf{问题：当前工作是不是在否定 \texttt{795}？}

\textbf{建议回答：} 不是。当前工作是在保留其最有价值实验经验的前提下，解决它在统一比较、模块化和持续迭代上的局限。

\item \textbf{问题：如果 \texttt{795} 这么有效，为什么还要重构？}

\textbf{建议回答：} 因为有效不等于可扩展。要把实验经验变成可训练、可验证、可持续迭代的系统，就必须做结构层面的重组。
\end{itemize}

\subsection{第 11 页：Bowman：参考路线的定位}
\begin{itemize}
\item \textbf{问题：Bowman 在这里到底指什么？是某篇论文，还是一个优化器？}

\textbf{建议回答：} 更准确地说，Bowman 在这里代表的是一类显式传播模型下的直接物理目标函数优化思路，而不只是某一个具体 CG 实现。

\item \textbf{问题：为什么要单独把 Bowman 拿出来讲？}

\textbf{建议回答：} 因为当前系统的最终物理收尾和可解释性，主要都建立在这条路线之上。如果不把 Bowman 讲清楚，后面的 hybrid 主线就容易被误解为纯经验拼接。
\end{itemize}

\subsection{第 12 页：Bowman：源码中的落点与当前吸收方式}
\begin{itemize}
\item \textbf{问题：当前代码里哪里最能代表 Bowman 的落点？}

\textbf{建议回答：} 最直接的是 \texttt{ai\_holography/hybrid.py} 中的 \texttt{bowman\_cg\_refine(...)}，以及 \texttt{compare\_bowman\_ai.py} 这种对比入口。

\item \textbf{问题：当前系统是不是把 Bowman 原样搬过来了？}

\textbf{建议回答：} 不是。当前系统保留了 Bowman 的物理优化内核，但前面增加了 candidate selection 和可微 refine，使它更多承担最终物理收尾而不是从任意初值艰难起步。
\end{itemize}

\subsection{第 13 页：Lin：参考路线的定位}
\begin{itemize}
\item \textbf{问题：Lin 路线最重要的思想是什么？}

\textbf{建议回答：} 最重要的是把问题改写到更适合学习的表示域，让网络先学位置域 amplitude / phase，再恢复 hologram 初值，而不是直接硬回归最终相位。

\item \textbf{问题：为什么这种改写有意义？}

\textbf{建议回答：} 因为这样通常训练更稳定，也更有利于学习与目标结构相关的初始化规律，并和后续物理优化更自然地衔接。
\end{itemize}

\subsection{第 14 页：Lin：源码中的落点与当前吸收方式}
\begin{itemize}
\item \textbf{问题：当前系统对 Lin 的吸收是不是只停留在网络结构？}

\textbf{建议回答：} 不是。除了位置域预测结构，当前训练还叠加了 teacher phase distillation 和 flat-top 质量损失，因此它已经被改造成适配当前任务的学习型初始化器。

\item \textbf{问题：评价 Lin checkpoint 最该看什么？}

\textbf{建议回答：} 不能只看 supervised loss，更要看它作为 init 时，能否帮助后续 hybrid polish 进入更优 basin。
\end{itemize}

\subsection{第 15 页：为什么当前系统最终采用 Bowman + Lin 的混合路线}
\begin{itemize}
\item \textbf{问题：为什么不是二选一，而是一定要混合？}

\textbf{建议回答：} 因为 Bowman 和 Lin 解决的是不同阶段的问题。Lin 更擅长学习“从哪里开始更好”，Bowman 更擅长解决“最后怎样收得更稳、更符合物理约束”。

\item \textbf{问题：这两条路线在当前系统里是竞争关系吗？}

\textbf{建议回答：} 不是，更准确地说是前后分工关系。当前有效性恰恰来自它们在不同阶段承担不同职责。
\end{itemize}

\subsection{第 16 页：四条路线的角色对照：795 / Bowman / Lin / 当前混合系统}
\begin{itemize}
\item \textbf{问题：如果老师问“你到底做的是哪一条路线”，怎么回答？}

\textbf{建议回答：} 最好的回答是：\texttt{795} 是实验先验和参考基线，Bowman 是最终物理收尾机制，Lin 是学习型初始化器，而当前真正交付的是把三者整合起来的混合系统。

\item \textbf{问题：这张表在答辩里最适合什么时候用？}

\textbf{建议回答：} 最适合在老师追问“各条路线关系”或“你工作的边界”时使用，它能把前面的信息快速收口成清晰对照。
\end{itemize}

\subsection{第 17 页：新程序的总体设计目标}
\begin{itemize}
\item \textbf{问题：新程序最核心的设计原则是什么？}

\textbf{建议回答：} 是三者整合：\texttt{795} 提供实验先验，Lin2025 提供学习型初始化，Bowman/hybrid 提供最终物理收尾。

\item \textbf{问题：为什么不只保留物理优化，而还要引入神经网络？}

\textbf{建议回答：} 因为纯物理优化对初值敏感，而学习型初始化器的价值恰恰在于帮助系统更稳定地进入高质量 basin。
\end{itemize}

\subsection{第 18 页：新程序整体运行流程}
\begin{itemize}
\item \textbf{问题：神经网络在当前系统中到底负责什么？}

\textbf{建议回答：} 主要负责初始化，而不是直接产出最终可交付解。最终质量仍由后端物理优化来保障。

\item \textbf{问题：为什么流程里要有 candidate selection？}

\textbf{建议回答：} 因为不同来源的初始化在不同任务下表现并不一致，统一选优能把“初值敏感”从经验问题变成可管理流程。
\end{itemize}

\subsection{第 19 页：新程序的目录结构}
\begin{itemize}
\item \textbf{问题：为什么要拆成这么多目录？}

\textbf{建议回答：} 因为训练、推理、物理优化、远程工作流和汇报整理承担的职责不同，拆分之后更便于维护、验证和继续演进。

\item \textbf{问题：这种拆分最大的实际收益是什么？}

\textbf{建议回答：} 最大收益是边界更清晰，后续做改动时能更容易判断问题来自目标定义、训练策略，还是物理优化与 benchmark 逻辑。
\end{itemize}

\subsection{第 20 页：版本演进总览}
\begin{itemize}
\item \textbf{问题：为什么版本史不能只看 notebook 标题？}

\textbf{建议回答：} 因为多个版本存在标题、说明文字和真实逻辑不同步的情况，必须结合 cell 内容、change log 和关键结构变化来重建。

\item \textbf{问题：最值得记住的关键版本有哪些？}

\textbf{建议回答：} 当前最关键的是 \texttt{v30 / v42a / v43 / v53 / v56}，而当前推荐交付版本是 \texttt{v58}。
\end{itemize}

\subsection{第 21 页：v2-v18 / v19-v39：从可运行到主线成形}
\begin{itemize}
\item \textbf{问题：为什么说 \texttt{v30} 是关键拐点？}

\textbf{建议回答：} 因为它是第一代真正稳定、可复用的主线骨架，后续很多版本本质上都还在围绕 \texttt{v30} 主方向做修正。

\item \textbf{问题：早期版本最大的贡献是什么？}

\textbf{建议回答：} 不是指标最优，而是把远程训练、teacher 池和任务定义这些基础设施逐步搭建起来。
\end{itemize}

\subsection{第 22 页：v40-v50：集成 notebook 与历史表阶段}
\begin{itemize}
\item \textbf{问题：为什么中期版本强调单变量验证？}

\textbf{建议回答：} 因为只有把关键变量拆开验证，才能分清哪些改善是有效主线，哪些只是偶然共振或过渡尝试。

\item \textbf{问题：\texttt{v42a} 为什么重要？}

\textbf{建议回答：} 因为它清楚表明 reference generation 中的 \texttt{efficiency\_floor} 本身就是决定结果的重要变量。
\end{itemize}

\subsection{第 23 页：v51-v58：高分辨率、结构一致性与当前推荐版本}
\begin{itemize}
\item \textbf{问题：为什么当前推荐的是 \texttt{v58}，不是 \texttt{v53} 或 \texttt{v56}？}

\textbf{建议回答：} 因为 \texttt{v58} 在保留后期有效主线的同时，把导出、summary 路径和版本一致性收口得更完整，更适合作为稳定交付版。

\item \textbf{问题：为什么说 \texttt{v58} 不是新的算法代际？}

\textbf{建议回答：} 因为它更多解决的是结构和一致性问题，而不是再引入一个决定算法方向的新机制。
\end{itemize}

\subsection{第 24 页：模块 1：目标与参考相位模块}
\begin{itemize}
\item \textbf{问题：这个模块的核心价值是什么？}

\textbf{建议回答：} 核心价值是统一目标定义，并把历史 \texttt{795} 参考相位变成当前系统可管理、可适配、可比较的先验来源。

\item \textbf{问题：为什么参考相位还值得长期保留？}

\textbf{建议回答：} 因为它承载了真实实验中被验证过的结构经验，是当前系统里非常重要的一类初始化来源。
\end{itemize}

\subsection{第 25 页：模块 2：传播与损失模块}
\begin{itemize}
\item \textbf{问题：为什么传播和损失模块可以看成“物理底盘”？}

\textbf{建议回答：} 因为训练、refine 和 hybrid polish 都共享这套传播建模和损失定义，它决定了全系统是否在同一物理与评价口径下运行。

\item \textbf{问题：如果没有这一层统一，会发生什么？}

\textbf{建议回答：} 不同路线之间就难以公平比较，benchmark 也会失去说服力，因为每条路线可能其实在优化不同的问题。
\end{itemize}

\subsection{第 26 页：模块 3：AIHolographyPipeline}
\begin{itemize}
\item \textbf{问题：这一层和 runner 的区别是什么？}

\textbf{建议回答：} pipeline 更偏流程编排和统一问题构造，runner 更偏 route 调度、候选选优和整条路线执行。

\item \textbf{问题：为什么这一层重要？}

\textbf{建议回答：} 因为它把原来分散的逻辑整理成可扩展的统一流程，后续增加新初始化器或 profile 会更容易。
\end{itemize}

\subsection{第 27 页：模块 4：HybridHolographyRunner}
\begin{itemize}
\item \textbf{问题：runner 解决的主要痛点是什么？}

\textbf{建议回答：} 主要是把不同来源初值纳入统一选优流程，显式缓解了物理优化对初值敏感的问题。

\item \textbf{问题：为什么不能手工指定一种初始化一直用下去？}

\textbf{建议回答：} 因为不同任务和不同 checkpoint 下，最佳起点并不总是同一种来源，统一选优比固定假设更稳。
\end{itemize}

\subsection{第 28 页：模块 5：Bowman-style hybrid polish}
\begin{itemize}
\item \textbf{问题：为什么还必须保留这一步物理收尾？}

\textbf{建议回答：} 因为这一步是当前系统结果稳定性、物理一致性和最终质量的核心保障，不能简单被学习模型替代。

\item \textbf{问题：当前工作和 Bowman 的关系应如何表述？}

\textbf{建议回答：} 最准确的表述是补强而非替代。我们保留 Bowman 的强项，同时用学习型初始化器来改善其对初值敏感的弱点。
\end{itemize}

\subsection{第 29 页：模块 6：Lin2025 数据集与 teacher 策略}
\begin{itemize}
\item \textbf{问题：为什么 teacher 不能只用 \texttt{795} 或只用 hybrid？}

\textbf{建议回答：} 因为两者各有价值。\texttt{795} 提供实验经验结构，hybrid 提供当前系统内部的高质量物理解，两者结合更有利于学习稳定而有效的初始化规则。

\item \textbf{问题：teacher 会固定不变吗？}

\textbf{建议回答：} 不会。当前系统的设计允许 teacher 随 benchmark 结果持续更新，这也是系统可持续迭代的重要原因。
\end{itemize}

\subsection{第 30 页：模块 7：Lin2025 网络结构}
\begin{itemize}
\item \textbf{问题：为什么网络不直接输出最终 hologram phase？}

\textbf{建议回答：} 因为直接回归最终相位通常更不稳定，也更难和后续物理优化结合。位置域预测更符合当前混合系统的角色分工。

\item \textbf{问题：这是不是意味着网络能力不够？}

\textbf{建议回答：} 不是。这里是有意进行系统分工，让网络专注于它最擅长的初始化任务，而不是强行承担全部求解职责。
\end{itemize}

\subsection{第 31 页：模块 8：Lin2025 训练循环}
\begin{itemize}
\item \textbf{问题：为什么训练里要放这么多损失项？}

\textbf{建议回答：} 因为目标不是单纯拟合 teacher，而是让网络产出对后续物理优化真正有帮助的初值，因此需要同时约束表征质量和任务质量。

\item \textbf{问题：多个 checkpoint 的意义是什么？}

\textbf{建议回答：} 不同 checkpoint 代表不同侧重点，后续 benchmark 可以帮助判断哪种训练目标更有利于最终全链路表现。
\end{itemize}

\subsection{第 32 页：模块 9：统一 benchmark}
\begin{itemize}
\item \textbf{问题：为什么 benchmark 是系统核心，而不只是结果展示？}

\textbf{建议回答：} 因为它决定了我们是否能在统一口径下回答“相对基线到底提升了多少”，没有 benchmark 就很难形成可复现的系统性结论。

\item \textbf{问题：当前 benchmark 最重要的作用是什么？}

\textbf{建议回答：} 一是比较 route，二是帮助筛选 checkpoint，三是把经验判断转化为可追踪的结果证据。
\end{itemize}

\subsection{第 33 页：模块 10：远程训练与本地分析}
\begin{itemize}
\item \textbf{问题：为什么要拆成远程训练和本地分析两段？}

\textbf{建议回答：} 因为 GPU 训练和统一分析的运行环境需求不同，拆开之后更稳定，也更有利于 checkpoint 管理和长期复现。

\item \textbf{问题：当前推荐 notebook 为什么是 \texttt{v58}？}

\textbf{建议回答：} 因为它在后期主线基础上把导出、summary 路径和结构一致性都收得更完整，适合作为当前交付入口。
\end{itemize}

\subsection{第 34 页：结果对比}
\begin{itemize}
\item \textbf{问题：这页最重要的一句话结论是什么？}

\textbf{建议回答：} 最重要的结论是，当前混合系统已经能在统一 benchmark 下稳定体现出相对于旧路线的系统性增益。

\item \textbf{问题：为什么不逐个数字展开讲？}

\textbf{建议回答：} 因为真正重要的是整体趋势和结构性判断，具体数值可以按评审关注点再展开，而不是一开始就陷入细节。
\end{itemize}

\subsection{第 35 页：相对于 795 的优势}
\begin{itemize}
\item \textbf{问题：相对于 \texttt{795} 的最大进步是什么？}

\textbf{建议回答：} 最大进步是从单一路线经验求解，转变为可训练、可比较、可复现的混合优化系统。

\item \textbf{问题：这种优势只体现在结果吗？}

\textbf{建议回答：} 不只体现在结果，还体现在方法论和工程组织方式上，例如统一 benchmark、checkpoint 选优和远程训练工作流。
\end{itemize}

\subsection{第 36 页：相对于 Bowman/hybrid 的价值}
\begin{itemize}
\item \textbf{问题：当前系统是不是已经“超过” Bowman 了？}

\textbf{建议回答：} 更准确的说法是，在部分结果上找到了优于纯物理基线的 Pareto 点，但系统本质上仍然建立在 Bowman 风格物理收尾之上。

\item \textbf{问题：为什么一定要强调“补强而非替代”？}

\textbf{建议回答：} 因为这能准确反映当前工作的创新点，也能避免把系统误解为纯 AI 替代传统方法。
\end{itemize}

\subsection{第 37 页：当前不足}
\begin{itemize}
\item \textbf{问题：当前最关键的短板是哪一个？}

\textbf{建议回答：} 如果从对实验闭环的影响看，最关键的是 line-flat 目标和真实实验链路尚未完全重新对齐；如果从工程稳定性看，则是 checkpoint 波动和高分辨率一致性仍需继续收敛。

\item \textbf{问题：为什么要主动讲不足？}

\textbf{建议回答：} 因为这有助于界定当前成果适用范围，也能说明后续工作方向是清楚的，而不是只展示局部最优结果。
\end{itemize}

\subsection{第 38 页：最终结论}
\begin{itemize}
\item \textbf{问题：如果老师只问一句“你这项工作的核心贡献是什么”，怎么回答？}

\textbf{建议回答：} 可以回答：核心贡献是把 \texttt{795} 的实验先验、Lin2025 的学习型初始化和 Bowman 的物理优化整合成一条可训练、可比较、可复现的完整主线。

\item \textbf{问题：下一步最重要的工作是什么？}

\textbf{建议回答：} 不是继续盲目堆版本，而是围绕 line-flat 目标、真实实验闭环、高分辨率一致性和 checkpoint 稳定性继续收敛。
\end{itemize}

\section{附录：Bowman 与 Lin 的极简答辩口径}

这一附录保留一份更短的复述版，用于老师在自由提问阶段突然追问“你怎么理解 Bowman 和 Lin”时快速回答。由于前文第 11--16 页已经完整展开，这里不再重复长篇解释，只保留三个最需要记住的结论：

\begin{enumerate}[label=\arabic*.]
\item \textbf{对 Bowman 的理解：}
不是把它狭义看成一个 CG 优化器，而是把它看成一条显式传播模型下、直接面向物理目标函数求解的路线。因此它最适合承担最终物理收尾。

\item \textbf{对 Lin 的理解：}
不是让网络直接生硬回归最终 hologram，而是把问题改写到更适合学习的表示域，让网络优先学习初始化规律。因此它最适合承担学习型初始化器角色。

\item \textbf{对当前系统的理解：}
真正有效的地方不在于二选一，而在于让 Lin 负责“从哪里开始更好”，让 Bowman 负责“最后怎么收得更稳、更符合物理约束”，再由 runner 和 benchmark 把两者组织成统一主线。
\end{enumerate}

如果需要用一段最短的正式口径来概括，可以直接采用如下表述：

\begin{quote}
我对 Bowman 的理解，是把它看成一条直接面向物理目标函数求解的路线，它的优势在于物理一致性强、最终收尾稳定。\\
我对 Lin 的理解，是把它看成一条学习型初始化路线，它的优势在于能够学习更好的起点。\\
当前系统的关键，不是让两者彼此替代，而是把 Lin 作为初始化器，把 Bowman 作为最终物理优化器，再通过统一 runner 和 benchmark 把它们接成一条完整主线。
\end{quote}

\section{收尾建议}

如果需要在正式汇报结束时使用一段较完整的总结陈述，可采用如下口径：

\begin{quote}
当前工作的主要价值，在于已经将原始 \texttt{795} 的实验经验、Lin2025 的学习型初始化以及 Bowman 的物理优化方法，整合为一套具有统一 benchmark、远程训练与本地分析能力的完整系统。现阶段这条主线已经证明了自身的有效性，后续工作的重点将不再是简单叠加版本，而是继续围绕 line-flat 目标、真实实验闭环、高分辨率一致性和 checkpoint 稳定性做进一步收敛。
\end{quote}

\end{document}
